% Discussion and conclusion (outline section 5): four run-in paragraphs, no new numbers.
\section{Discussion and conclusion}\label{sec:discussion}

Frozen features keep what the pretraining labels separate, and merging quark flavor is the costliest
single merge. On the pretraining classes a linear probe separates the \Kfull{} classes less
well at every coarser level of the class hierarchy (figure~\ref{fig:f2a}(a)), and the step that removes
quark flavor costs the most in wrongly ordered pairs of jets (table~\ref{tab:probes}). % CHECK: MlstCmpJcTwoCoarseOverMid > MlstCmpJcTwoMidOverRes > MlstCmpJcTwoResOverFull > 1
Merging the QCD classes costs only the QCD classes. Random partitions into \Kcoarse{} classes, built to
cut across the physical groups, lose less than the physical \Kcoarse-class label set, so the class count alone does
not set the cost. Per class, decays merged with their lepton-flavor counterparts pay the most, while
decays merged with their quark-flavor counterparts carry the largest summed cost
(figure~\ref{fig:f3}; data release, \texttt{jetclass2\_per\_class.csv}). % CHECK: MlstPerClassCoarseTopClass is a lepton-flavor class; MlstPerClassCoarseQuarkShare > MlstPerClassCoarseLeptonShare
A single flavor boundary placed at random among the four-quark decays keeps flavor readable even for
pairs it does not cut (figure~\ref{fig:f4}). Fine-tuning shrinks the cost but does not remove it where every input is present: on JetClass the
\Kcoarse-class model gives \MlstFtJcOneErrCoarseOverFullBoot{} than the \Kfull-class model, while the
\Kres- and \Kmid-class models are consistent with it (table~\ref{tab:jetclass}).
% CHECK: MlstFtJcOneErrCoarseOverFullBoot passes 3 sigma; MlstFtJcOneErr{Res,Mid}OverFullBoot read "consistent"
On top tagging and quark/gluon discrimination, which lack some inputs, the fine-tuned label sets show no
resolved difference (tables~\ref{tab:benchtop} and~\ref{tab:benchqg}).
With the missing inputs set to zero, a linear probe on the frozen features of the \Kfull-class model
matches the randomly initialized network on top tagging (\MlstLpTopRthirtyFullOverInitBnZeroPr{}); with the
substitution chosen on validation it reaches \MlstLpTopRthirtyFullOverInitBnPr{} of that network, and
\MlstLpQgRthirtyFullOverInitBnPr{} on quark/gluon discrimination (tables~\ref{tab:fill_main}
and~\ref{tab:fill}). On these benchmarks the zero value, more than the missing information, keeps the
frozen features from doing better than a random network. % CHECK: ...FullOverInitBnZeroPr reads "equal within"; ...FullOverInitBn above 1 by more than 2 SD on both tasks
The two effects act on one part of the input. Flavor reaches the network through the displacement and
particle-type inputs, because b and c hadrons travel a few millimeters before they decay and leptons are
identified by type. These are the inputs the benchmarks lack, so a label-set difference that lives in
flavor likely has little to act on there.

The results translate into concrete choices for a pretraining label set. Keep every distinction that a
downstream discriminant sums over, because a merge removes the discriminant from the output layer
whatever the frozen features keep: Sophon's $X\to b\bar{b}$ discriminant cannot be formed from the
\Kmid-class labels, and the $D_{bc}$ discriminant of the boosted $|V_{cb}|$ measurement exists only at
\Kfull{} and \Kres{} classes (table~\ref{tab:usecase}). Keep the flavor boundaries: on the pretraining classes a single one lowers the error of the linear probe,
and on JetClass the cost after fine-tuning appears only at the merge of quark flavor
(figures~\ref{fig:f2a}--\ref{fig:f4} and table~\ref{tab:jetclass}). Keep the QCD classes apart where a downstream task separates QCD jets by
flavor: merging them leaves the resonant classes unchanged but costs the QCD pairs
(table~\ref{tab:pairs}) and the one discriminant that sums QCD classes
(table~\ref{tab:usecase}). % CHECK: MlstCmpJcTwoQcdResOverFull > 1 and the QCD pairs of tab:pairs show 162 above 188
Add an auxiliary mass regression: it puts the jet mass into the frozen features (table~\ref{tab:mass} and
figure~\ref{fig:f9a}). Its effect on mass sculpting is resolved only for the \Kres-class model at a QCD
efficiency of \SculptEffHigh, where it reduces it (table~\ref{tab:sculptboot}). Its cost to the classification of the pretraining classes
is small; after fine-tuning on top tagging neither label set with it shows a change beyond the uncertainty
of $R_{30}$, and \MlstMassNRunsWord{} runs per model cannot exclude a modest one
(section~\ref{sec:mass}). % CHECK: MlstCmpMlpResMassOverRes < 1; the sculpting clause is gallery_placement.F10B_STATED's 162mass-162 at eps_B 10% (the only mass-regression difference beyond 3 sigma in sculpting_boot.json, PI 2026-11-09); as section 4.6 states, MlstFtTopRthirtyResMassOverResPr and MlstFtTopRthirtyCoarseMassOverCoarsePr read "equal within" under the Methods two-run rule (raw R30 today: 1.12 and 0.89; 0.90 and 0.92); if either states a direction beyond MlstBootTestMaxRthirtyFt, rewrite both sections
For benchmarks with missing inputs, compare label sets by fine-tuning with the missing input features
set to zero, which gives every label set identical inputs; input dropout during pretraining, which
replaces some inputs with zeros, is the published alternative~\cite{omnilearn,jpjepa}. The label set matters most where the features stay frozen or the fine-tuning set is small: after
fine-tuning on JetClass the \Kcoarse-class labels give \MlstFtJcOneErrCoarseOverFullBootEFour{} at
\FtSizeEFour{} jets and \MlstFtJcOneErrCoarseOverFullBoot{} at \FtSizeESix{}.
% CHECK: MlstFtJcOneErrCoarseOverFullBootEFour waits for the 1e4 bootstrap after the equal-stopping read-out (2026-11-09: about 7% more errors, z = 4.6 with the binomial term); top and quark/gluon at small training sets wait for the equal-stopping results
Report the variation between pretraining runs, because the
label-set effects here are measured against it (section~\ref{sec:uncertainties}) and a single
pretrained model cannot tell the two apart.

The study has four limitations. First, all pretraining and every evaluation but one use a single
simulation family; the Herwig quark/gluon sample changes the generator for one task only. Second, one
architecture and one model size are studied. Third, each level of the class hierarchy has
\MlstNRunsWord{} pretraining runs and every other configuration fewer. % the outline's wording "fewer for the controls" is widened: the mass models (two runs) and the kinematic-only model (one run) are not controls
Fourth, the networks are initialized with a different scheme from Sophon's code, so their absolute
performance need not match Sophon's published model; every comparison is between models trained here.

Which distinctions a pretraining label set keeps, more than how many classes it has, shapes which
discriminants an analysis can still form and how well frozen features separate decays by flavor. A
per-particle reconstruction objective recovers jet flavor without flavor labels: frozen features of
machine-learned particle flow separate $b$ jets from light and $c$ jets~\cite{mlpf_fm}, whereas jet-level
classification keeps only what its label set distinguishes. A
label set built for transfer should therefore keep the classes that analyses sum into discriminants
and its flavor boundaries, and an auxiliary target such as the jet mass can add to the frozen features a
quantity that no classification label rewards, for little loss in classification.
